\section{Dataset Description}
The project uses the TrashNet dataset (resized version)~\cite{trashnet}, a public dataset of photographs of single waste items placed on a white background under varying lighting and orientation. The original dataset has six classes (glass, paper, cardboard, plastic, metal, trash). For this project it was filtered to \textbf{four material classes}: glass, plastic, paper and metal.

\tablestyle
\begin{table}[h]
\centering
\caption{Classes used from TrashNet}
\begin{tabularx}{\textwidth}{|l|r|X|}\hline
\hdrrow \hdr{Class} & \hdr{Approx.\ images} & \hdr{Remarks} \\
\hline
Glass   & 501 & Bottles, jars \\
\hline
Paper   & 594 & Sheets, printed paper \\
\hline
Metal   & 410 & Cans, tins \\
\hline
Plastic & 482 & Bottles, containers \\
\hline
\textbf{Total} & $\mathbf{\approx 1{,}987}$ & Filtered from 2,527 original images \\
\hline
\end{tabularx}
\end{table}

\subsection{Train / Validation / Test Split}
The filtered dataset was split into training (\textbf{70\%}), validation (\textbf{15\%}) and test (\textbf{15\%}) sets using the split-folders library with a fixed random seed (\textbf{42}) for reproducibility. The test set is never used during training and gives an unbiased estimate of real-world performance.
